\chapter[Persistence and explicit resolutions]{Two-parameter persistence, Schur syzygies and surface-order coefficients}
\label{pers:chapter}
A persistence diagram records how objects and maps change under an ordered
family of inclusions. A projective resolution records the generators,
relations and relations among relations of one module. These descriptions
meet naturally over a polynomial parameter ring. They must still be
distinguished from the infinite locally gentle algebra attached to the
surface: the parameter ring describes the filtration, not a replacement
surface algebra.

The three supplied papers motivate these interfaces
\cite{PersSurfaceI,PersSurfaceII,PersSurvey}. We use their surface-order,
lattice and finite-presentation ideas with the corrections already proved
in this book. In particular, a matrix satisfying consecutive-product zero
is a complex, not automatically a resolution; Fitting base change is not
exact base change; and the existence of a generic resolution is not a
finite terminating algorithm. The worked example below proves its own
exactness and exhibits an actual topological realization.

\section{The polynomial ring encodes the two commuting directions}
Let $k$ be a field and give $\mathbb N^2$ its coordinatewise order. A
persistence module is a covariant functor $P:\mathbb N^2\to\operatorname{Vect}_k$.
Writing $P_{a,b}=P(a,b)$, put
\[
 R=k[x,y],\quad \deg x=(1,0),\quad\deg y=(0,1),\qquad
 \mathcal P=\bigoplus_{a,b\ge0}P_{a,b}.
\]
Multiplication by $x$ and $y$ is the map to the next coordinate in the
corresponding direction. Functoriality makes their square commute.
Conversely the graded pieces and multiplication maps of an
$\mathbb N^2$-graded $R$-module define such a functor. General transition
maps are multiplication by $x^{c-a}y^{d-b}$ when $(a,b)\le(c,d)$.
The exponents are nonnegative in this direction, a useful sign check.

This gives an equivalence for the stated discrete diagrams. Finite
generation is an additional condition, not a consequence of having
finite-dimensional pieces. A module can have infinitely many independent
births. Finite cellular input with finitely many birth grades gives the
finite presentations used here. A real-parameter filtration needs an
appropriate discretization or a different coefficient algebra before
being treated as a module over $k[x,y]$.

We write $R(-a,-b)$ for a free module whose generator has grade $(a,b)$.
Over a field, a finite graded module has a minimal graded free resolution
of length at most two. Minimality means that every differential entry
belongs to $(x,y)$; tensoring with $k=R/(x,y)$ then kills the differentials.
The numbers and grades of the surviving generators are the multigraded
Betti numbers. The field hypothesis matters: $\mathbb Z[x,y]$ has
homological dimension three, and the regular sequence $(2,x,y)$ supplies
a length-three example. The particular resolution below also works
integrally, but that does not make every integral resolution length two.

\section{A complete three-grade example}
Set
\[
 M=R/(x^2,xy,y^2).
\]
Its basis is $1,x,y$, in grades $(0,0),(1,0),(0,1)$. All other graded
pieces vanish. The two nonzero proper transition maps send $1$ to $x$
and $1$ to $y$. In particular, for comparable grades $\nu\le\mu$,
\[
 \rho_M(\nu,\mu)=\operatorname{rank}(M_\nu\to M_\mu)
 =\begin{cases}1&\nu,\mu\in\{(0,0),(1,0),(0,1)\},\\0&\text{otherwise.}\end{cases}
\]
This computes this module's rank invariant; rank invariants do not in
general classify all two-parameter modules.

\begin{figure}[htbp]\centering
\begin{tikzpicture}[>=Stealth,x=1.35cm,y=1.15cm]
\foreach \a in {0,1,2,3}{\foreach \b in {0,1,2,3}{
 \draw[gray!60] (\a,\b) circle (2pt);}}
\foreach \a/\b in {0/0,1/0,0/1}{\fill[teal!70!black] (\a,\b) circle (4pt);}
\draw[->,very thick,teal!70!black] (.09,0)--(.88,0);
\draw[->,very thick,teal!70!black] (0,.09)--(0,.88);
\draw[->,gray] (-.3,-.35)--(3.5,-.35) node[right]{$a$};
\draw[->,gray] (-.35,-.3)--(-.35,3.5) node[above]{$b$};
\node[below] at (0,-.35) {$0$};\node[below] at (1,-.35) {$1$};
\node[left] at (-.35,1) {$1$};
\node[above right] at (0,0) {$1$};\node[above] at (1,0) {$x$};
\node[right] at (0,1) {$y$};
\node[anchor=west,align=left] at (3.9,1.7) {Filled point: one copy of $k$\\Open point: zero\\Only the two thick arrows survive};
\end{tikzpicture}
\caption{The actual rank diagram of $R/(x,y)^2$. Moving right or up outside
its three-point support kills the class. Incomparable grades are not joined
by a persistence map.}\label{pers:rankfigure}
\end{figure}

\begin{theorem}[A minimal bigraded resolution]\label{pers:resolution}
The following complex is a minimal graded free resolution of $M$:
\[
 0\longrightarrow F_2\xrightarrow{d_2}F_1\xrightarrow{d_1}R
       \longrightarrow M\longrightarrow0,
\]
where
\[
 \begin{split}
 F_1&=R(-2,0)\oplus R(-1,-1)\oplus R(0,-2),\\
 F_2&=R(-2,-1)\oplus R(-1,-2),\qquad
 d_1=\begin{pmatrix}x^2&xy&y^2\end{pmatrix},\quad
 d_2=\begin{pmatrix}-y&0\\x&-y\\0&x\end{pmatrix}.
 \end{split}
\]
\end{theorem}
\begin{proof}
Each column of $d_2$ has the indicated source grade, and multiplication
gives $d_1d_2=0$. If $d_2(a,b)=0$, its first and last entries give
$ya=xb=0$. Since $R$ is a domain, $a=b=0$.

Suppose $x^2f+xyg+y^2h=0$. Reducing modulo $x$ shows $x\mid h$,
say $h=xb$. Dividing the relation by $x$ and reducing modulo $y$
shows $y\mid f$, so write $f=-ya$. Substitution now gives
$y(-xa+g+yb)=0$, hence $g=xa-yb$. Thus
$(f,g,h)=d_2(a,b)$. This proves exactness at $F_1$. The image of
$d_1$ is its displayed ideal, giving exactness at $R$ and the
augmentation. Every entry of both matrices has positive grade, so the
resolution is minimal. The divisibility argument also works in
$\mathbb Z[x,y]$, where $x,y$ remain prime and nonzerodivisors.
\end{proof}

The multigraded Betti numbers are consequently
\[
 \beta_{0,(0,0)}=1,\quad
 \beta_{1,(2,0)}=\beta_{1,(1,1)}=\beta_{1,(0,2)}=1,\quad
 \beta_{2,(2,1)}=\beta_{2,(1,2)}=1.
\]
They are dimensions of the corresponding graded $\operatorname{Tor}^R_i(M,k)$.
Taking degreewise Euler characteristics gives the same Hilbert series by
two independent descriptions:
\[
 H_M(s,t)=1+s+t
 =\frac{1-s^2-st-t^2+s^2t+st^2}{(1-s)(1-t)}.
\]
After forgetting the separate grading this becomes
$(1-3z^2+2z^3)/(1-z)^2=1+2z$. The numerator records the alternating
resolution shifts, not a characteristic polynomial whose zeros locate
arithmetic $L$-zeros.

One can also see the dual calculation directly. Ignoring shifts only to
write the matrices, the dual differential is
\[
 d_2^*=\begin{pmatrix}-y&x&0\\0&-y&x\end{pmatrix}.
\]
The equations $-yA+xB=-yB+xC=0$ imply
$(A,B,C)=(x^2q,xyq,y^2q)$, again using coprimality. Thus
$\operatorname{Ext}^0_R(M,R)=\operatorname{Ext}^1_R(M,R)=0$.
The second Ext is generated by $e_1,e_2$ with relations
$ye_1=0$, $xe_1=ye_2$ and $xe_2=0$; its $k$-basis is
$e_1,e_2,xe_1$. This keeps the syzygy calculation distinct from merely
counting the support points of $M$.

\section{A bifiltered CW realization and the syzygies it creates}
Start with one vertex and an oriented loop $e$, both born at grade $(0,0)$.
Attach three two-cells by degree-one maps to that loop, at grades
\[
 D_{20}:(2,0),\qquad D_{11}:(1,1),\qquad D_{02}:(0,2).
\]
At grade $(a,b)$ include precisely the cells whose birth grades are at
most $(a,b)$. Every boundary is already present, so these are actual
subcomplex inclusions and their squares commute. The persistence cellular
chain modules are
\[
 C_2=F_1,\qquad C_1=R e,\qquad C_0=R,\qquad
 \partial_2=(x^2\ \ xy\ \ y^2),\quad\partial_1=0.
\]
Therefore $H_1=M$ exactly. No assertion that every abstract persistence
module has this particular geometric form is needed.

The loop survives before any disk is present and dies when the first disk
appears. Two disks attached along their common boundary form a sphere.
The first two persistent $H_2$ generators are the compatible cell cycles
\[
 -yD_{20}+xD_{11}\quad\text{at }(2,1),\qquad
 -yD_{11}+xD_{02}\quad\text{at }(1,2).
\]
Theorem~\ref{pers:resolution} proves
$H_2\cong R(-2,-1)\oplus R(-1,-2)$; there are no three-cells to kill
them. At $(2,2)$ all three disks give two independent sphere classes.
Thus the same matrix $d_2$ is a relation among the killing cells and a
positive persistent homology class in the larger CW complex.
This auxiliary CW space need not be a Riemann surface: three disks along
one circle are not a manifold. It is a topological realization of the
filtration module, not another claimed cellular dessin.

\section{Schur modules organize the generators and the first syzygies}
Assume now $\operatorname{char}k=0$, put $V=kx\oplus ky$, and let
$\mathrm{GL}(V)$ act on $R=\operatorname{Sym}V$ by linear substitution.
The ideal $(x,y)^2$ is invariant. Its degree-two generators form
$\operatorname{Sym}^2V$, with basis $x^2,xy,y^2$.
The linear syzygies form the kernel of multiplication
\[
 V\otimes\operatorname{Sym}^2V\longrightarrow\operatorname{Sym}^3V.
\]
There is a natural equivariant map
\[
 \delta:\bigwedge^2V\otimes V\longrightarrow V\otimes\operatorname{Sym}^2V,
 \qquad (v\wedge w)\otimes z\longmapsto
       v\otimes wz-w\otimes vz.
\]
Multiplication kills its image. On the two basis vectors its values are
exactly the two columns of $d_2$:
\[
 (x\wedge y)\otimes x\mapsto x\otimes xy-y\otimes x^2,
 \qquad
 (x\wedge y)\otimes y\mapsto x\otimes y^2-y\otimes xy.
\]
They are independent, and the kernel has dimension $6-4=2$, so $\delta$
identifies it with $\det V\otimes V=S_{(2,1)}V$. The Littlewood--Richardson
rule gives
\[
 V\otimes\operatorname{Sym}^2V
 \cong\operatorname{Sym}^3V\oplus S_{(2,1)}V.
\]
Here the partition convention uses symmetric rows. The equivariant total-
graded resolution has multiplicity spaces
$k,\operatorname{Sym}^2V,S_{(2,1)}V$ in generator degrees $0,2,3$.
The diagonal torus recovers their bigraded weights
$(2,0),(1,1),(0,2)$ and $(2,1),(1,2)$.

The grading distinction is essential. General $\mathrm{GL}_2$ substitutions
mix $x$ and $y$, so they preserve total degree, not the two original
filtration coordinates. Only the diagonal torus preserves each bigrading.
In positive characteristic the displayed maps and our exact resolution
remain valid, but semisimple direct-sum representation theory must not
be assumed. Nor does applying a Schur functor termwise to an arbitrary
exact sequence preserve exactness, as the earlier counterexamples show.

\section{Fitting data, exactness and specialization}
For a finite presentation $R^p\xrightarrow D R^q\to N\to0$ define
\[
 \operatorname{Fitt}_j(N)=I_{q-j}(D),\qquad
 I_0(D)=R,
\]
where $I_r$ is the ideal of $r$-minors, $I_r=R$ for $r\le0$, and
$I_r=0$ for $r>\min(p,q)$. The ideal is independent of the chosen finite
presentation. Its formation commutes with arbitrary base change because
minors are polynomials in the presentation entries. For our cyclic module,
\[
 \operatorname{Fitt}_0(M)=(x,y)^2,\qquad
 \operatorname{Fitt}_j(M)=R\quad(j\ge1).
\]
The support is the origin, but its Fitting scheme retains a length-three
thickening. This is not the same datum as the rank invariant of the
ordered graded pieces.

The $2$-minors of $d_2$ are $y^2,-xy,x^2$, so their ideal is also
$(x,y)^2$. Expected ranks are $1,2$ and this ideal has grade two;
$x^2,y^2$ is a regular sequence. These facts agree with the finite-free
acyclicity criterion of Buchsbaum--Eisenbud \cite{PersBE}, independently
of our direct proof. Over a general Noetherian ring, a height count cannot
silently replace the criterion's grade/regular-sequence hypotheses.
The criterion concerns a finite complex over a commutative ring; it does
not certify every infinite path-algebra resolution.

Arbitrary specialization does not preserve that exactness. At $x=y=0$
both matrices vanish, leaving positive homology $k^3,k^2$ in degrees
one and two. More informatively, after setting only $y=0$, the complex
has
\[
 H_0=k[x]/(x^2),\qquad H_1=(k[x]/(x))^2,\qquad H_2=0.
\]
These follow immediately from $d_1=(x^2,0,0)$ and the two surviving
entries $x$ of $d_2$. They are the derived-base-change Tor groups.
Flat base change, suitable Tor vanishing, or a new exactness proof is
needed to obtain an ordinary specialized resolution. A generic parameter
space and a small differential residual do not remove those conditions.

\section{A correct finite space of complexes and its incidence model}
For variety language in this section take $k$ algebraically closed.
Let $V_i$ have dimensions $n_i$ for $i=0,1,2$. Fix feasible ranks
$a,b$ with
\[
 a\le\min(n_0,n_1),\qquad b\le\min(n_1,n_2),\qquad a+b\le n_1.
\]
The reduced variety
\[
 \mathcal C_{a,b}=\{(A,B):V_2\xrightarrow B V_1\xrightarrow A V_0,
       \ AB=0,\ \operatorname{rank}A\le a,\ \operatorname{rank}B\le b\}
\]
parametrizes based complexes with these upper bounds, not isomorphism
classes of resolutions. Consider flags
$L\subset K\subset V_1$ with $\dim L=b$, $\dim K=n_1-a$, and impose
\[
 \operatorname{im}B\subset L\subset K\subset\ker A.
\]
Over the flag variety the allowed maps form the vector bundle
\[
 \operatorname{Hom}(V_1/K,V_0)\oplus\operatorname{Hom}(V_2,L).
\]
Its total space is smooth and irreducible. Projection is proper onto
$\mathcal C_{a,b}$: for any complex with the rank bounds, first extend
$\operatorname{im}B$ to a $b$-plane inside $\ker A$, then extend it to an
$(n_1-a)$-plane there. On the full-rank locus the choices are uniquely
$L=\operatorname{im}B$, $K=\ker A$. Full ranks exist by feasibility,
so this projection is birational. This proves the asserted incidence
resolution of the reduced rank-bound variety.

For $(n_0,n_1,n_2)=(1,3,2)$ and $(a,b)=(1,2)$, the two flag spaces
coincide in a plane, and the base is $\mathbf{Gr}(2,k^3)$. The vector
bundle has fiber dimension $1+4$ and base dimension two, giving total
dimension seven. Full-rank points are exact in positive degrees because
$\operatorname{im}B=\ker A$ and $B$ is injective. The zero pair is also a
point of the complex variety and is not exact. Our polynomial matrices
map the punctured plane to the exact locus; the origin maps to that
nonexact boundary. A specialization of the universal matrices therefore
still needs its own homological check.

Bruns' generic-resolution construction uses a directed limit of rings
which kills positive homology in the limit \cite[Theorem 1, pp.35--36]{PersBruns}.
It does not assert that some finite stage always becomes exact, nor
that every specialization of the resulting matrices preserves exactness.
The finite example in this chapter makes its stronger local claims by
explicit calculation rather than by invoking an unspecified generic object.

\section{Coefficients in a genuine infinite dessin order}
Return to the actual nodal interval dessin and distinguish its variables:
\[
 A=k[\xi,\eta]/(\xi\eta),\qquad
 \Lambda=k[[\xi,\eta]]/(\xi\eta),\qquad u=\xi+\eta.
\]
Both $A$ and $\Lambda$ are infinite dimensional over $k$.
The completed order is free of rank two over $k[[u]]$, on $1,\xi$,
with $\xi^2=u\xi$ and $\eta=u-\xi$. Multiplication by $u$ is injective,
as is clear in its two normalization branches. Set $E=\Lambda/u^2\Lambda$.
It has basis $1,\xi,u,u\xi$ over $k$, nontrivial arrow actions, and
\[
 0\longrightarrow\Lambda\xrightarrow{u^2}\Lambda\longrightarrow E
       \longrightarrow0
\]
is a genuine free resolution. The module $E$ is finite length, not a
$k[[u]]$-free lattice. Its finite dimension does not turn its parent order
into a finite-dimensional substitute.

The family $E\otimes_k M$ is a persistence module valued in continuous
$\Lambda$-modules: each nonzero piece is $E$, and its structure maps are
the identity or zero as in Figure~\ref{pers:rankfigure}. The actions of
$x,y$ commute with the fixed action of $u,\xi,\eta$. There is no
identification of either persistence direction with the order uniformizer.

More precisely put $\mathcal B=\Lambda\otimes_k k[x,y]$. Tensor the
length-one $\Lambda$-resolution just displayed with
Theorem~\ref{pers:resolution} over $k$. Its finite total complex has
free $\mathcal B$-ranks
\[
 (\operatorname{rank}F_0,\operatorname{rank}F_1,
   \operatorname{rank}F_2,\operatorname{rank}F_3)=(1,4,5,2)
\]
and resolves $E\otimes_k M$. If $Q_i$ denotes the order resolution and
$P_j$ the persistence resolution, the total differential on $Q_i\otimes P_j$
is
\[
 D=\partial_Q\otimes1+(-1)^i1\otimes d_P.
\]
Order the degree-$n$ terms as $P_n\oplus(Q_1\otimes P_{n-1})$.
Suppressing only the tensor symbols, the three matrices are
\[
 D_1=(d_1\ \ u^2),\qquad
 D_2=\begin{pmatrix}d_2&u^2I_3\\0&-d_1\end{pmatrix},\qquad
 D_3=\begin{pmatrix}u^2I_2\\-d_2\end{pmatrix}.
\]
The cross terms cancel, so consecutive products are zero. For exactness,
write a positive-degree cycle as $(a,b)$, with
$d a+u^2b=0$ and $d b=0$. Modulo $u^2$ the first component is a cycle
in $E\otimes_k P_\bullet$. That complex is exact in positive degrees,
since tensoring vector spaces over $k$ is exact. Hence
$\bar a=d\bar c$. Lift $\bar c$ to $c$ in the free term and write
$a-dc=u^2e$. The original equation becomes $u^2(de+b)=0$.
Multiplication by $u^2$ is injective on every free term, so $b=-de$.
Consequently $(a,b)=D(c,e)$, as required. The same argument includes
the top degree by putting missing terms equal to zero. Degree-zero
homology is $\mathcal B/(u^2,x^2,xy,y^2)=E\otimes_k M$.
This proves exactness without assuming that an arbitrary tensor product
of resolutions stays exact. The ranks follow from
$(1+z)(1+3z+2z^2)=1+4z+5z^2+2z^3$.

This construction is compatible with the original center action. It does
not shorten the infinite resolution of the nodal residue simple, make
$E$ a lattice, or construct Frobenius. Restriction to the commutative
coefficient ring $k[[u]][x,y]$ doubles these free ranks. Determinant
functors apply to the specified finite perfect complex, with virtual
rank $1-4+5-2=0$, but determinant multiplicity and arithmetic comparison
remain additional questions.

\section{Quantum data, monotonicity and the next usable interface}
An arbitrary time-varying interaction graph is not a bifiltration: an
edge can disappear at a later time. One valid alternative fixes finite
birth pairs $(a_e,b_e)$ for edges and includes an edge exactly when both
coordinates are reached, then takes the flag complex. Another uses a
cumulative union over times up to the first coordinate and decreasing
interaction thresholds in the second. Both give inclusions; the second
records cumulative history rather than individual snapshots. Vertices
and higher cells must have compatible birth grades. Zigzag persistence
is a different appropriate tool when deletions must be retained.

Quantum channels are linear on density matrices, but two unrelated
channels need not commute. Assigning them to the two directions does
not define a persistence functor unless the required squares commute.
One may instead extract a declared classical correlation graph, choose
a monotone filtration, and compute its homology with the conventions
above. The data-to-distance rule and its statistical uncertainty still
need justification.

For the research program this yields concrete trainable and checkable
targets: propose birth grades, presentation entries or a low-degree
syzygy; test inclusions, gradings, the matrix identities and actual
exactness; record Betti and Fitting data before interpreting them.
A loss penalizing $D^2$ checks only the complex equation. The exact
module example, the flag-incidence model and the surface-order coefficient
resolution show how to add the missing checks. None identifies those
finite invariants with a completed arithmetic $L$-function, provides a
positive arithmetic spectral limit, or proves GRH.

\section{The executable staircase and quantum laboratory}
The canonical on-page tool accompanies this chapter; other chapters link
here. Proposition~\ref{thm:staircase} generalizes the worked square-ideal
resolution to every strict Artinian staircase. Four displayed ideals have
lengths $3,6,13,16$, with custom natural exponents also allowed. The tool
shows generator and adjacent-syzygy degrees, standard monomials, exact
comparable-grade maps and cellular incidence ranks. Coefficients in $E$
multiply field dimensions by four; the displayed free ranks over
$\Lambda[x,y]$ instead come from the $u^2$ cone. These are different rings.

Specializing the free resolution at a point over $\mathbb F_2$ computes
its fiber-complex homology. At the origin its dimensions are
$(1,m,m-1)$; at each of the three nonzero $\mathbb F_2$ points they vanish,
since a pure-power generator becomes a unit. This is derived fiber data,
not the CW homology at a filtration grade. The signed maximal minors and
Hilbert numerator are checked independently with symbolic polynomials.

The quantum tab uses the single saved trained four-qubit QCNN unitary $U$
and the channel $T(\rho)=\operatorname{Tr}_{2,3}(U\rho U^*)$.
The designed inputs have their last three qubits in $|000\rangle$.
Their first qubit runs over eight equatorial states
$(|0\rangle+e^{2\pi i j/8}|1\rangle)/\sqrt2$ and both poles.
This is an out-of-training-distribution illustrative family, not new
training or a held-out performance claim. Birth labels are fixed from the
input population $(I+Z_0)/2$: south $0$, equator $500$, north $1000$.
They are retained at the output; recomputing them from output energy
would not justify the same identity inclusion.

Floating-point input/output trace distances are scaled by $1000$ and
rounded upward. Their frozen integer graph labels are saved with a SHA-256
checksum and checkpoint provenance. The inequality of output and input
labels is verified directly, never imposed by taking their minimum.
Proposition~\ref{thm:channel-filtration} explains the exact-channel contractivity;
the numerical records are not interval certificates. For the recorded labels,
all finite boundary matrices and induced ranks are computed exactly over
$\mathbb F_2$. At grade $(400,500)$ the input contains an equatorial
8-cycle and an isolated south pole, with $(\beta_0,\beta_1)=(2,1)$.
Its $H_1$ map to $(600,500)$ has rank one, whereas its map to $(800,1000)$
has rank zero. The plotted two-parameter rank table fixes the target and
varies comparable source grades; it is not a complete multiparameter barcode.
Rotating the drawing changes the input-qubit Bloch layout only. The same
layout is used at output and does not embed its two-qubit state space.

Independent row elimination and the browser's bit-column implementation
agree on both Betti numbers and inclusion ranks for the saved cases.
These executable finite examples complement the proofs and the cited
persistence theory. The worked square-ideal module now has a kernel-checked ungraded
free resolution, described below. General staircase exactness, graded and
categorical interfaces, the coefficient cone and channel norm contractivity
remain separate formalization tasks. No learned
component or finite homology calculation supplies the unresolved positive
arithmetic determinant bridge.

\section{What the polynomial Lean resolution actually verifies}
The file \texttt{PersistenceResolution.lean} uses
$R=\operatorname{MvPolynomial}(\mathrm{Fin}\,2,k)$ for an arbitrary field.
Its free coordinate modules are $F_1=R\times R\times R$ and
$F_2=R\times R$, with the matrices of Theorem~\ref{pers:resolution}.
The augmentation is the actual $R$-linear quotient map to
$R/(x^2,xy,y^2)$. The proof of \texttt{d1\_kernel} quantifies over
all polynomial coefficient vectors; it is not a sampled scalar-rank test.

The reusable lemma \texttt{relation\_decomposition} starts with a
commutative domain and a pair $x,y$ such that $x\ne0$ and
$x\mid yz$ implies $x\mid z$. From
$x^2a+xyb+y^2c=0$ it constructs $p,q$ with
$(a,b,c)=(-yp,xp-yq,xq)$. For the actual polynomial variables, the
regularity premise is discharged by viewing $R$ as polynomials in $x$
with coefficients in $k[y]$. A polynomial is divisible by $x$ exactly
when its constant coefficient vanishes; multiplication by the nonzero
$y$ cannot kill such a coefficient. Thus the final kernel theorem has
no unproved regularity hypothesis. The first and last coordinates then
prove \texttt{d2\_injective}; \texttt{d1\_range} identifies the image
with the specified ideal. Augmentation exactness and surjectivity complete
the sequence, and the free terms are proved projective.

The field $k$ carries the residue action
$r\cdot a=\operatorname{constantCoeff}(r)a$.
The free modules retain ordinary polynomial multiplication. Both actions
are explicit in the source, with action identities checked in the kernel.
For every $R$-linear map from the appropriate free term to this residue
field, precomposition with $d_1$ or $d_2$ is zero, since all displayed
entries contain $x$ or $y$. These are actual Hom statements with their
module actions, not merely entrywise evaluations. Abstract Ext-group
identifications and the graded minimal-resolution interface are not yet
packaged as formal theorems.

All theorem types and axiom dependencies are printed. Only the standard
Lean foundations occur; no project axioms or incomplete proofs are used.
This establishes the entire underlying ungraded worked resolution and its
residue Hom differential calculation. It does not formalize every staircase,
the CW model, Schur equivariance, the $E$-coefficient tensor cone or any
arithmetic determinant comparison. Those distinctions remain visible on
the local verification page beside the executable laboratory.
